\documentclass[11.5pt,a4paper]{book}
\usepackage[utf8]{inputenc}
\usepackage[T1]{fontenc}
\usepackage[french]{babel}
\usepackage{geometry}
\geometry{papersize={148mm,210mm}, top=2.5cm, bottom=2.5cm, left=2cm, right=2cm}

\usepackage{helvet}
\renewcommand{\familydefault}{\sfdefault}
\usepackage{xcolor}
\usepackage{tcolorbox}
\tcbuselibrary{skins, breakable}

\definecolor{CyrulnikBlue}{HTML}{212F3D}
\definecolor{GreyText}{HTML}{566573}

\newtcolorbox{NoteClinique}{
  colback=white,
  colframe=CyrulnikBlue,
  fonttitle=\bfseries,
  title=Anamnèse Clinique,
  arc=0mm,
  breakable
}

\begin{document}

\chapter{Le Cerveau Primitif et le Rapt de la Parole}

\textit{Celui qui n'ose pas dire est un homme qui marche dans le brouillard de son propre corps. Ce n'est pas qu'il n'ait rien à dire ; c'est que la porte est verrouillée de l'intérieur, par une clé forgée très tôt, dans la forge de la peur et du conformisme.}

\section{Le Guetteur Invisible}
Au plus profond de nos replis cérébraux, là où la raison n'a pas encore de prise, se trouve l'amygdale. Pour le patient inhibé, l'amygdale n'est plus une sentinelle ; elle est devenue un geôlier. Lorsque nous sommes confrontés à une autorité, ou simplement à l'idée de déplaire, ce petit noyau déclenche ce que les neurosciences appellent le \textit{Hijack} émotionnel. 

L'information court-circuite le cortex préfrontal — le siège de notre pensée choisie et de notre liberté. Elle file par une "voie basse", plus rapide de quelques millisecondes, directement vers les centres de la survie. Le résultat est immédiat : la bouche s'assèche, les mots se fragmentent, et cette sensation de "grumeau" dans la gorge apparaît. Ce n'est pas votre volonté qui flanche, c'est votre biologie qui a choisi la sidération.

\begin{NoteClinique}
\textbf{Le cas de Marc, 42 ans.} \\
Cadre supérieur, brillant dans ses analyses écrites, Marc s'effondre dès qu'il doit demander une augmentation ou contester une directive absurde. Il décrit une sensation de "disparition physique". L'IRMf révélerait chez lui ce que nous pressentons : une hypo-activation du cortex préfrontal ventromédian. Marc n'a pas "perdu ses mots", il a perdu, à l'instant T, la capacité neuronale d'évaluer sa propre valeur.
\end{NoteClinique}

\section{La Chimie du Renoncement}
Chaque fois que nous nous taisons alors que notre âme hurle, nous payons un tribut hormonal. Le "non-dit" n'est pas un acte neutre. Le cerveau, percevant le danger social, inonde le système de cortisol — l'hormone du stress. Mais contrairement au zèbre qui fuit le lion, l'homme inhibé ne dépense pas cette énergie. Elle stagne.

À long terme, cet excès de cortisol devient neurotoxique. Il sature l'hippocampe, cette zone de la mémoire et de l'apprentissage. Voilà pourquoi, après une conversation où nous n'avons pas osé nous affirmer, nous nous sentons hébétés, incapables de nous souvenir des arguments que nous aurions dû employer. Le silence a littéralement "anesthésié" notre intelligence.

\section{La Théorie Polyvagale : Au-delà du "Combat-Fuite"}
Pour comprendre le silence de Léo ou de Marc, il ne suffit pas de parler de stress. Il faut invoquer les travaux du Dr Stephen Porges sur le nerf vague. Notre système nerveux autonome possède deux branches, mais le nerf vague se divise lui-même en deux circuits distincts : le complexe vagal ventral (la sécurité sociale) et le complexe vagal dorsal (l'extinction).

\begin{NoteClinique}
\textbf{L'extinction dorsale.} \\
Lorsque l'affirmation de soi est perçue comme un risque d'exclusion ou de mort symbolique, le cerveau ne choisit pas la fuite, mais le "collapsus". C'est le complexe vagal dorsal qui prend les commandes. La tension artérielle chute, l'esprit s'embrume, et le sujet entre dans un état de soumission archaïque. On ne "peut" plus parler, car physiologiquement, le corps s'est mis en mode survie par l'effacement.
\end{NoteClinique}

\section{La Mémoire du Corps : L'Enfance comme Laboratoire}
Le cerveau est une éponge plastique qui se sculpte au gré des expériences. L'inhibition de l'adulte prend racine dans les "petites morts" de l'enfance. Chaque fois qu'une expression de soi a été réprimée par une autorité (parentale, scolaire) au profit du "ton calme" ou de la "discrétion", le vmPFC a appris à se désactiver pour protéger l'individu du rejet. 

Nous voyons en IRM que ces circuits de l'inhibition sont devenus des autoroutes neuronales. S'affirmer à 40 ans demande alors non pas de la "volonté", mais un véritable dé-câblage de ces sentiers de la peur appris très tôt.

\section{Le Coût du Silence : Une Neurotoxicité Silencieuse}
Chaque renoncement est une micro-agression neurologique. L'accumulation de cortisol non métabolisé par l'action (puisque l'action d'affirmation est bloquée) altère la substance blanche et diminue la densité du cortex préfrontal. 

Le but de ce livre est de proposer un protocole de "résonance" pour rallumer ces zones éteintes, pour que Marc puisse à nouveau habiter son vmPFC et que le "grumeau" dans la gorge se dissolve dans une parole juste et libérée.

\end{document}
